% GENERATED FILE. Edit canonical lessons, then run npm run generate:books.
% canonical-source-hash: fnv1a64:a60d793b75e8265a
% canonical-lessons: UR-C12-pani, UR-C12-chai, UR-C12-doodh, UR-C12-roti

\chapter{Water, Tea, Milk, Bread}
\label{ch:ur-food-words}
\hlchaptermodality{12}

\begin{chapteropening}
\textbf{By the end of this chapter:} \emph{I can name water, tea, milk and bread in Urdu, say what I like using the frame Chapter 8 built, and sort each word as inherited or Persian.}

Say mujhe roṭī pasand hai, reusing Chapter 8's dative-experiencer frame on a plain noun instead of an infinitive, and sort pānī/dūdh/roṭī as inherited against chāy's Persian loan — with pānī reusing lenā's verb frame and dūdh correcting a tempting false cognate with English dough.
\end{chapteropening}

\section[pānī]{\ur{پانی} — \textquotedblleft{}water,\textquotedblright{} and \ur{لینا} put back to work}

\label{lesson:UR-C12-pani}



\begin{quote}
Say \textbf{merā dil} — my heart, closing the body words, and recall that its PIE root reached Urdu by the Persian road while the same root reached Sanskrit's \emph{hṛdaya} by a different one. This chapter turns to the four things you would actually ask for. \textbf{\ur{لینا}}, \textquotedblleft{}to take,\textquotedblright{} from three chapters back, gets its first real use since it was taught.
\end{quote}

\subsection*{You'll want to know first — one word}
\begin{quote}
\textbf{\ur{پانی}} — \emph{pānī} — \textbf{water}
\end{quote}

\textbf{\ur{پانی}} is masculine.

\begin{grammarlens}[title={Grammar Lens: naming what you'd take}]
\begin{quote}
\textbf{\ur{میں} \ur{پانی} \ur{لیتا} \ur{ہوں۔}} — \emph{maiṅ pānī letā hūṅ.} — I take/get water (man speaking). \textbf{\ur{میں} \ur{پانی} \ur{لیتی} \ur{ہوں۔}} — \emph{maiṅ pānī letī hūṅ.} — I take/get water (woman speaking).
\end{quote}

\textbf{\ur{لینا}}'s stem \textbf{le-} slots a noun in front of it the same way any transitive verb does. You already hold the whole verb; \textbf{\ur{پانی}} is the first noun you have had to put in front of it.
\end{grammarlens}

\begin{cousinweb}
\textbf{pānī} continues Sanskrit \emph{pānīya}, \textquotedblleft{}water, drink,\textquotedblright{} built on \textbf{pā-}, \textquotedblleft{}to drink,\textquotedblright{} from Proto-Indo-European \textbf{*peh\textsubscript{3}-}. English \textbf{water} is unrelated — it comes from an entirely different PIE root. But \emph{*peh\textsubscript{3}-} itself did reach English, by a different door: Latin \emph{potare}, \textquotedblleft{}to drink,\textquotedblright{} gave English \textbf{potion} and, by way of a sense narrowed to \textquotedblleft{}a poisoned drink,\textquotedblright{} \textbf{poison}. \textbf{\ur{پانی}} is, underneath, \textquotedblleft{}the drinkable thing\textquotedblright{} — the same idea English kept for what you drink, not for the substance itself.
\end{cousinweb}

\subsection*{Your turn}
\begin{itemize}
\raggedright
  \item \emph{Say it:} \textbf{pānī} — water; then \textbf{merā pānī}
  \item \emph{Say it:} \textbf{maiṅ pānī letā hūṅ}, then \textbf{maiṅ pānī letī hūṅ}
  \item \emph{Connect:} \textbf{pānī} $\leftarrow$ \emph{pā-} $\to$ Latin \emph{potare} $\to$ English \textbf{potion}, \textbf{poison} — not \textbf{water}
\end{itemize}

\begin{tcolorbox}[breakable,colback=teal!4,colframe=teal!35!black,title={Before you move on}]
Is \textbf{\ur{پانی}} related to English \textbf{water}? (\textbf{No.}) Name two English words that do share its root. (\textbf{Potion} and \textbf{poison}.) Say \textquotedblleft{}I take/get water,\textquotedblright{} once as a man would, once as a woman would.

Sources: \href{https://en.wiktionary.org/wiki/\%D9\%BE\%D8\%A7\%D9\%86\%DB\%8C}{Wiktionary: \ur{پانی}}.
\end{tcolorbox}
\section[chāy]{\ur{چائے} — \textquotedblleft{}tea,\textquotedblright{} which walked to Urdu rather than sailed}

\label{lesson:UR-C12-chai}



\begin{quote}
Say \textbf{merā pānī} — my water, underneath just \textquotedblleft{}the drinkable thing.\textquotedblright{} Here is the drink you would politely ask for right after it, and a letter from Chapter 9 makes a second appearance.
\end{quote}

\subsection*{You'll want to know first — one word}
\begin{quote}
\textbf{\ur{چائے}} — \emph{chāy} — \textbf{tea}
\end{quote}

\begin{quote}
\textbf{\ur{میری} \ur{چائے۔}} — \emph{merī chāy.} — \textquotedblleft{}my tea.\textquotedblright{}
\end{quote}

\textbf{\ur{چائے}} is feminine.

\subsection*{The letters in this word}
From the right edge: \textbf{\ur{چ}} \emph{ch}, already yours from \emph{sochnā} and \emph{samajhnā}; then \textbf{\ur{ا}} long \emph{ā}; then \textbf{\ur{ئ}}, hamza-ye, the vowel-glide letter \textbf{\ur{بھائی}} introduced; then \textbf{\ur{ے}}, broad final \emph{e} from \emph{kaise}. \textbf{\ur{چائے}}'s \textbf{\ur{ئ}} does the same job it did in \textbf{\ur{بھائی}}: keeping \textbf{ā} and the vowel after it from colliding.

\begin{cousinweb}
\textbf{\ur{چائے}} is borrowed from Classical Persian \textbf{\ur{چای}} (\emph{chāy}), which traces back to Northern Chinese \textbf{chá}, carried overland through Central Asia into Persian before Urdu ever had it. That is the same road \textbf{\ur{خدا}} and \textbf{\ur{دل}} travelled in earlier chapters — Persian into Urdu — just with a Chinese starting point this time instead of one further west.

Tea that instead reached Europe by \textbf{sea} kept a different form, from southern coastal Chinese, which is why English says \textbf{tea} and French says \emph{thé} instead of \emph{chai}. One drink, two words, sorted by whether it travelled by land or by water.
\end{cousinweb}

\subsection*{Your turn}
\begin{itemize}
\raggedright
  \item \emph{Say it:} \textbf{chāy} — tea; then \textbf{merī chāy}
  \item \emph{Name:} the hamza-ye letter's two words so far — \textbf{bhāī}, \textbf{chāy}
  \item \emph{Say it:} the road — \textbf{chá … chāy … chāy}
  \item \emph{Contrast:} \textbf{merā pānī} — \textbf{merī chāy}
\end{itemize}

\begin{tcolorbox}[breakable,colback=teal!4,colframe=teal!35!black,title={Before you move on}]
Which three languages did \textbf{\ur{چائے}} pass through, in order? (\textbf{Chinese, then Persian, then Urdu.}) Why does English say \textquotedblleft{}tea\textquotedblright{} instead of \textquotedblleft{}chai\textquotedblright{}? (\textbf{It arrived by sea, carrying a different Chinese form.}) Which earlier Urdu word already used the hamza-ye letter? (\textbf{\ur{بھائی}}.)

Sources: \href{https://en.wiktionary.org/wiki/\%DA\%86\%D8\%A7\%D8\%A6\%DB\%92}{Wiktionary: \ur{چائے}}.
\end{tcolorbox}
\section[dūdh]{\ur{دودھ} — \textquotedblleft{}milk,\textquotedblright{} and a false friend worth naming before you fall for it}

\label{lesson:UR-C12-doodh}



\begin{quote}
Say \textbf{merī chāy} — my tea, which walked overland from China through Persian before reaching Urdu. This word is what turns it milky, and it hands the two-eyed \textbf{\ur{ھ}} a sixth letter to aspirate.
\end{quote}

\subsection*{You'll want to know first — one word}
\begin{quote}
\textbf{\ur{دودھ}} — \emph{dūdh} — \textbf{milk}
\end{quote}

\begin{quote}
\textbf{\ur{میرا} \ur{دودھ۔}} — \emph{merā dūdh.} — \textquotedblleft{}my milk.\textquotedblright{}
\end{quote}

\textbf{\ur{دودھ}} is masculine.

\subsection*{The letters in this word}
From the right edge: \textbf{\ur{د}} \emph{d}, then \textbf{\ur{و}} carrying long \emph{ū}, then \textbf{\ur{د}} again, then \textbf{\ur{ھ}}. \textbf{\ur{دھ}} extends the two-eyed \emph{he} to a sixth base letter: it has now aspirated \textbf{\ur{ج،} \ur{ٹ،} \ur{ڑ،} \ur{چ،} \ur{ک}}, and here it aspirates \textbf{\ur{د}}, giving \emph{dh}.

\begin{cousinweb}
\textbf{dūdh} continues Sanskrit \emph{dugdhá}, literally \textquotedblleft{}milked,\textquotedblright{} from Proto-Indo-European \textbf{*d\textsuperscript{h}ewg\textsuperscript{h}-}, \textquotedblleft{}to be fit, to yield, to produce.\textquotedblright{} That root also gives Sanskrit \textbf{duh-}, \textquotedblleft{}to milk\textquotedblright{} — the verb sitting directly under this noun.

Here is the trap. English \textbf{dough} looks like the obvious match, and it is \textbf{not}: \emph{dough} comes from a completely different PIE root, \textbf{*d\textsuperscript{h}eig\textsuperscript{h}-}, \textquotedblleft{}to knead, to form\textquotedblright{} — the root behind the shape of bread, not the substance of milk. The real English cousin of \textbf{\ur{دودھ}} is a quieter word: \textbf{doughty}, \textquotedblleft{}brave, capable,\textquotedblright{} from the very same \textbf{*d\textsuperscript{h}ewg\textsuperscript{h}-}, \textquotedblleft{}fit for, of use.\textquotedblright{} A doughty person and milked milk share an ancestor; bread dough does not.
\end{cousinweb}

\subsection*{Your turn}
\begin{itemize}
\raggedright
  \item \emph{Say it:} \textbf{dūdh} — milk; then \textbf{merā dūdh}
  \item \emph{Count:} six letters now carrying \textbf{\ur{ھ}} — \textbf{\ur{ج،} \ur{ٹ،} \ur{ڑ،} \ur{چ،} \ur{ک،} \ur{د}}
  \item \emph{Reject:} \textbf{\ur{دودھ}} is NOT related to English \textbf{dough}
  \item \emph{Connect:} \textbf{\ur{دودھ}} $\leftarrow$ \emph{*d\textsuperscript{h}ewg\textsuperscript{h}-} $\to$ English \textbf{doughty}, the real cousin
\end{itemize}

\begin{tcolorbox}[breakable,colback=teal!4,colframe=teal!35!black,title={Before you move on}]
Which English word looks like a match for \textbf{\ur{دودھ}} but is not one? (\textbf{Dough.}) Which English word actually is the cousin? (\textbf{Doughty.}) What Sanskrit verb sits directly under \textbf{\ur{دودھ}}? (\textbf{duh-, \textquotedblleft{}to milk.\textquotedblright{}})

Sources: \href{https://en.wiktionary.org/wiki/\%D8\%AF\%D9\%88\%D8\%AF\%DA\%BE}{Wiktionary: \ur{دودھ}}; \href{https://www.etymonline.com/word/doughty}{Etymonline: doughty}.
\end{tcolorbox}
\section[roṭī]{\ur{روٹی} — \textquotedblleft{}bread,\textquotedblright{} and \ur{پسند} closes the chapter}

\label{lesson:UR-C12-roti}



\begin{quote}
Three drinks so far: \textbf{pānī, chāy, dūdh} — and recall \textbf{\ur{دودھ}}'s real English cousin was \textbf{doughty}, not the tempting \textbf{dough}. This chapter's fourth and final word is the one you would eat, not drink.
\end{quote}

\subsection*{You'll want to know first — one word}
\begin{quote}
\textbf{\ur{روٹی}} — \emph{roṭī} — \textbf{bread}
\end{quote}

\begin{quote}
\textbf{\ur{میری} \ur{روٹی۔}} — \emph{merī roṭī.} — \textquotedblleft{}my bread.\textquotedblright{}
\end{quote}

\textbf{\ur{روٹی}} is feminine.

\begin{cousinweb}
\textbf{roṭī} continues Sauraseni Prakrit \emph{roṭṭa-}, from Sanskrit \emph{roṭikā}, \textquotedblleft{}a kind of bread.\textquotedblright{} That word is built inside Sanskrit itself, on a root meaning \textquotedblleft{}to \textbf{strike, to crush}\textquotedblright{} — the pounding and rolling that turns grain into flatbread. The trail stops there: no proposed Proto-Indo-European ancestor, no English cousin. \textbf{\ur{کان}} and \textbf{\ur{پڑھنا}} hit the same honest wall earlier in this book. Not every word gets a Greek or Latin relative, and saying so plainly is better than inventing one.
\end{cousinweb}

\begin{grammarlens}[title={Grammar Lens: \ur{مجھے} ... \ur{پسند} \ur{ہے}, one more time}]
\begin{quote}
\textbf{\ur{مجھے} \ur{روٹی} \ur{پسند} \ur{ہے۔}} — \emph{mujhe roṭī pasand hai.} — \textbf{I like bread.}
\end{quote}

The dative-experiencer frame from Chapter 8 was built for \textbf{-nā} infinitives — \emph{paṛhnā}, \emph{likhnā}. It works just as well for a plain noun: whatever pleases you sits in the subject slot, and \textbf{\ur{ہے}} agrees with it, not with you.
\end{grammarlens}

\subsection*{Your turn}
\begin{itemize}
\raggedright
  \item \emph{Say it:} \textbf{roṭī} — bread; then \textbf{merī roṭī}
  \item \emph{Say it:} \textbf{mujhe roṭī pasand hai}, then swap in \textbf{pānī}, then \textbf{dūdh}
  \item \emph{Sort:} inherited — \textbf{pānī, dūdh, roṭī}; Persian — \textbf{chāy}
  \item \emph{Say it:} all four foods with \textbf{merā} or \textbf{merī} correctly chosen
\end{itemize}

\begin{tcolorbox}[breakable,colback=teal!4,colframe=teal!35!black,title={Before you move on}]
Does \textbf{\ur{روٹی}} have a secure English cousin? (\textbf{No — the trail stops at Sanskrit.}) Of this chapter's four foods, which one is the Persian loan? (\textbf{\ur{چائے}}.) Say \textquotedblleft{}I like bread\textquotedblright{} using the frame from Chapter 8.

Sources: \href{https://en.wiktionary.org/wiki/\%D8\%B1\%D9\%88\%D9\%B9\%DB\%8C}{Wiktionary: \ur{روٹی}}.
\end{tcolorbox}
